\uuid{WdFd}
\titre{DL usuels/Calcul de limites}
\auteur{hanani ; mzari}
\organisation{exo7}
\licence{CC BY-NC-SA 4.0}
\datecreate{2020-01-20}
\niveau{L1}
\module{Analyse}
\chapitre{Développement limité}
\sousChapitre{Applications}
\difficulte{1}
\theme{Applications des DL, DL usuels, Calcul de limites}
\exo7qcmid{lille-2.developpements-limites.031}

\contenu{
\texte{
Soit $a\in \Rr$ et $f$ une fonction. On dit que $f(x)$ est un $o\left((x-a)^n\right)$ au voisinage de $a$ et on écrit $f(x)=o((x-a)^n)$ si $\displaystyle \lim _{x\to a}\frac{f(x)}{(x-a)^n}=0$. On note par $DL_n(a)f(x)$ le développement limité de $f$ à l'ordre $n$ en $a$. On note aussi $DL_n(a^+)f(x)$ (resp. $DL_n(a^-)f(x)$) le développement limité de $f$ à l'ordre $n$ à droite en $a$ (resp. à gauche en $a$).
}

\begin{question}
\qclassification{125.03}
\qtags{motcle=DL usuels/Calcul de limites}
Soit $\displaystyle f(x)=\frac{\ln x}{x^2-1}$ et $\displaystyle g(x)=\frac{\ln (1+x)}{2x+x^2}$. Parmi les affirmations suivantes, cocher celles qui sont vraies :
\begin{answers}
\good{$\displaystyle DL_1(0)g(x)=\frac{1}{2}-\frac{x}{2}+o(x)$.}
\bad{$f(1+x)=g(x)$ et $\displaystyle DL_1(1)f(x)=\frac{1}{2}-\frac{x}{2}+o(x)$.}
\bad{$\displaystyle \lim _{x\to 1}f(x)$ n'existe pas.}
\good{$\displaystyle \lim _{x\to 1}f(x)=\frac{1}{2}$.}
\end{answers}
\begin{explanations}
Les développements à l'ordre $2$ en $0$ de $\ln (1+x)$ donne :
$$\displaystyle DL_1(0)g(x)=\frac{1}{2}-\frac{x}{2}+o(x) \Rightarrow DL_1(1)f(x)=\frac{1}{2}-\frac{x-1}{2}+o(x-1)$$
car $f(x)=g(x-1)$. Enfin, $\displaystyle \lim _{x\to 1}f(x)=\lim _{x\to 0}g(x)=\frac{1}{2}$.
\end{explanations}
\end{question}
}
